\section{Entrevistas con representantes de usuarios finales}
\label{sec:entrevistas}

La nota final del enunciado establece textualmente:
\textit{``Se espera que en la definición formal de su proyecto (PC02)
se tomen en cuenta estos y otros resultados, incluyendo resultado de
entrevistas con representantes de sus usuarios finales''}.

Esta sección presenta la \textbf{guía operativa} elaborada por el
equipo a partir del prompt 06 (Gemini, versionado en
\texttt{prompts/respuestas/gemini\_06\_entrevistas.md}) y deja
explícitas las acciones que el equipo ejecutará para completar el
trabajo con datos reales.

\begin{cajaPrompt}[title={Estado actual --- pendiente de ejecución de campo}]
La guía operativa (perfiles, criterios de selección, bloques de
preguntas, métricas a registrar y plantilla de consentimiento) está
lista. Falta el trabajo de campo: contactar entrevistados, ejecutar
las sesiones, transcribir y consolidar hallazgos. La consolidación
final se incorporará en este mismo archivo en una iteración posterior
del informe.
\end{cajaPrompt}

\subsection{Perfil 1 --- Estudiante de últimos ciclos (\sk{1})}

\paragraph{Objetivo} Identificar puntos de dolor en la transición
académica-laboral y validar la aceptación de una interfaz de
entrenamiento basada en IA.

\paragraph{Criterios de selección}
\begin{itemize}
  \item Estudiante matriculado en 8\textsuperscript{vo},
        9\textsuperscript{no} o 10\textsuperscript{mo} ciclo de una
        carrera de ingeniería.
  \item Sin experiencia previa en prácticas preprofesionales o con
        máximo una experiencia técnica.
  \item Búsqueda activa de empleo declarada para los próximos 6 meses.
  \item Usuario recurrente de herramientas de videoconferencia
        (Zoom, Meet, Teams).
\end{itemize}

\paragraph{Bloques de preguntas}
\begin{itemize}
  \item \textbf{Contexto y experiencia previa}:
        \begin{itemize}
          \item ¿Cómo ha sido su proceso de preparación para las
                primeras postulaciones laborales?
          \item ¿Qué aspectos de una entrevista presencial o virtual le
                generan mayor incertidumbre?
        \end{itemize}
  \item \textbf{Expectativas frente a la IA}:
        \begin{itemize}
          \item ¿Qué valor esperaría obtener de un evaluador que no
                es humano?
          \item ¿Cómo reaccionaría si el nivel de dificultad de las
                preguntas cambia según sus respuestas anteriores?
        \end{itemize}
  \item \textbf{Metáfora del dojo y el aura}:
        \begin{itemize}
          \item Al imaginar un entorno de ``entrenamiento constante''
                tipo dojo, ¿qué elementos visuales consideraría
                motivadores?
          \item ¿De qué manera interpretaría un halo de luz que cambia
                de color mientras usted habla?
        \end{itemize}
  \item \textbf{Privacidad y ética}:
        \begin{itemize}
          \item ¿Qué condiciones exigiría para permitir que un
                algoritmo analice sus gestos faciales?
          \item ¿Cómo preferiría que el sistema le informe sobre el
                uso de sus datos tras finalizar la sesión?
        \end{itemize}
\end{itemize}

\paragraph{Métricas a registrar}
\begin{itemize}
  \item Frecuencia de mención de términos relacionados con
        ``ansiedad'' o ``nerviosismo''.
  \item Tiempo de respuesta ante la propuesta del aura
        (latencia en la reacción).
  \item Recuento de adjetivos positivos/negativos asociados al
        feedback automatizado.
\end{itemize}

\subsection{Perfil 2 --- Recién egresado (\sk{2})}

\paragraph{Objetivo} Evaluar la utilidad de simulaciones de alta
complejidad y la efectividad del plan de mejora personalizado para la
inserción laboral real.

\paragraph{Criterios de selección}
\begin{itemize}
  \item Egresado universitario con menos de 12 meses de haber culminado
        estudios.
  \item Postulante activo con al menos 3 entrevistas realizadas en el
        último trimestre.
  \item Acceso a equipo de cómputo propio con webcam y micrófono
        funcional.
\end{itemize}

\paragraph{Bloques de preguntas}
\begin{itemize}
  \item Experiencia en procesos reales: motivos de rechazo frecuentes,
        retroalimentación recibida.
  \item Funcionalidad del plan de mejora: información indispensable,
        uso cotidiano del plan.
  \item Metáfora y gamificación: rangos/ligas, sala virtual.
  \item Privacidad y datos: almacenamiento por la universidad,
        confianza en análisis de lenguaje corporal por IA.
\end{itemize}

\paragraph{Métricas a registrar}
\begin{itemize}
  \item Sugerencias específicas para el plan de mejora.
  \item Grado de interés en gamificación (escala 1--5 observada).
  \item Palabras clave recurrentes (``empleabilidad'', ``competencia'').
\end{itemize}

\subsection{Perfil 3 --- Coach de carrera o reclutador (\sk{8}, \sk{9})}

\paragraph{Objetivo} Contrastar las métricas automáticas del sistema
con los estándares profesionales de evaluación de talento.

\paragraph{Criterios de selección}
\begin{itemize}
  \item Profesional con más de 3 años de experiencia en reclutamiento
        IT o coaching laboral.
  \item Experiencia previa evaluando perfiles \textit{junior} o
        practicantes de ingeniería.
  \item Conocimiento básico de IA aplicada a RRHH.
\end{itemize}

\paragraph{Bloques de preguntas}
\begin{itemize}
  \item Criterios de evaluación: indicadores no verbales,
        progresión de dificultad.
  \item Validación de la herramienta: fiabilidad del análisis de
        muletillas y ritmo, sesgos.
  \item Interacción humano-IA: reacción ante respuestas incompletas,
        valor del peer mock.
  \item Privacidad institucional: garantías de seguridad, reportes
        agregados.
\end{itemize}

\paragraph{Métricas a registrar}
\begin{itemize}
  \item Términos técnicos de reclutamiento omitidos por la propuesta.
  \item Gestos de escepticismo o validación ante las métricas de
        lenguaje corporal.
  \item Tiempo dedicado a discutir la ética de la IA.
\end{itemize}

\subsection{Perfiles 4--7 --- Usuarios con discapacidad (\sk{4}--\sk{7})}

\paragraph{Objetivo} Validar la accesibilidad de la interfaz multimodal
y la eficacia de los canales alternativos de feedback.

\paragraph{Criterios de selección}
\begin{itemize}
  \item \sk{4} (visual): usuario dependiente de lector de pantalla o
        con baja visión severa.
  \item \sk{5} (auditiva): usuario con pérdida auditiva que utilice
        subtitulado o lengua de señas.
  \item \sk{6} (motora): usuario con movilidad reducida en manos que
        utilice software de asistencia.
  \item \sk{7} (cognitiva/neurodivergente): usuario diagnosticado
        dentro del espectro autista o con TDAH.
\end{itemize}

\paragraph{Bloques de preguntas (comunes y específicos)}
\begin{itemize}
  \item \textbf{Barreras de accesibilidad}: obstáculos en plataformas
        de entrevista por video actuales; interacción típica con
        asistentes de voz o comandos verbales.
  \item \textbf{Interfaz no convencional} (específico \sk{4}/\sk{5}):
        \begin{itemize}
          \item (\sk{4}) ¿Cómo esperaría recibir la información del
                aura si esta es puramente visual?
          \item (\sk{5}) ¿Qué tan útil le resulta la descripción
                textual del tono de voz del entrevistador?
        \end{itemize}
  \item \textbf{Carga cognitiva y ayuda} (específico \sk{6}/\sk{7}):
        \begin{itemize}
          \item (\sk{6}) ¿Qué dificultades prevé al interactuar con
                una interfaz que requiere respuestas en tiempo real?
          \item (\sk{7}) ¿De qué manera el avatar del entrevistador
                podría reducir la ansiedad social en lugar de
                aumentarla?
        \end{itemize}
  \item \textbf{Manejo de sensores}:
        \begin{itemize}
          \item ¿Qué inquietudes le genera que el sistema analice sus
                patrones de movimiento o voz específicos?
          \item ¿Cómo prefiere ser notificado si el sistema no logra
                procesar correctamente su imagen o audio?
        \end{itemize}
\end{itemize}

\paragraph{Métricas a registrar}
\begin{itemize}
  \item Identificación de fallos potenciales en el flujo de navegación
        mediante voz o teclado.
  \item Menciones explícitas al estándar WCAG o necesidades de
        personalización.
  \item Signos de fatiga cognitiva durante la descripción de
        funcionalidades complejas.
\end{itemize}

\subsection{Plantilla de consentimiento informado}

\begin{cajaCriterio}
\textbf{Título del proyecto}: Simulador de Entrevistas Laborales
Adaptativo (CC451 --- UNI). \\
\textbf{Investigadores}: Davila Santos, Serrano Arostegui, Poma
Navarro.

\paragraph{Propósito} Se le solicita participar en una entrevista para
recoger requisitos y percepciones sobre una plataforma de simulación
de entrevistas con IA. La información ayudará a mejorar la
accesibilidad y eficacia del sistema.

\paragraph{Procedimiento} La sesión durará entre 45 y 60 minutos. Se
registrará audio y se tomarán notas de la transcripción para análisis
académico exclusivo del curso CC451.

\paragraph{Confidencialidad} Sus datos personales serán anonimizados
mediante códigos. Las grabaciones se eliminarán al finalizar el ciclo
académico 2026-I. No se compartirá video crudo ni audio con entidades
externas al equipo de desarrollo y la cátedra.

\paragraph{Derechos} Su participación es voluntaria. Puede negarse a
responder cualquier pregunta o retirarse del proceso en cualquier
momento sin penalización alguna.

\paragraph{Consentimiento}
Yo, \rule{6cm}{0.4pt}, identificado con DNI/Código \rule{2.5cm}{0.4pt},
declaro haber sido informado sobre los objetivos del estudio y acepto
voluntariamente participar, autorizando la grabación de esta sesión
únicamente para fines de transcripción y análisis académico. \\[0.5cm]
Firma: \rule{5cm}{0.4pt} \quad Fecha: \rule{0.6cm}{0.4pt}/\rule{0.6cm}{0.4pt}/2026
\end{cajaCriterio}

\subsection{Plan de ejecución y consolidación}

\paragraph{Acciones del equipo}
\begin{enumerate}
  \item Reclutar al menos un entrevistado por cada perfil
        (\sk{1}, \sk{2}, \sk{8} o \sk{9}, y un representante por cada
        \sk{4}--\sk{7}).
  \item Realizar las entrevistas (45--60 min cada una) con
        consentimiento informado firmado.
  \item Transcribir las grabaciones y registrar hallazgos por
        entrevista.
  \item Consolidar patrones recurrentes, discrepancias y ajustes al
        \textit{backlog} de RF/RNF o casos de uso, con su justificación
        basada en los hallazgos.
  \item Incorporar la consolidación en este mismo archivo en la
        siguiente iteración del informe, con la siguiente estructura
        por entrevista: perfil del entrevistado, fecha y duración,
        modalidad, hallazgos clave por bloque temático, citas
        literales relevantes (con consentimiento explícito) y cambios
        al diseño del proyecto derivados, con referencia al RF/RNF o
        caso de uso afectado.
\end{enumerate}

\paragraph{Trazabilidad entrevistas a requerimientos} Por cada
hallazgo accionable se registrará: ID de la entrevista (E-NN),
hallazgo en una oración, requerimiento o caso de uso impactado
(\rf{N}, \rnf{N}, \cu{N}), tipo de impacto
(confirma / agrega / modifica / descarta).
